% ==============================================================================
% FILOSOFÍA DEL DERECHO - ACTIVIDAD 4 / SEMANA 6
% Estudio de caso: Poder, Estado y Derecho
% Archivo no autocontenido: requiere la plantilla base del flujo de trabajo.
% ==============================================================================

\documentclass[
spanish,
letterpaper, oneside
]{article}

% INFORMACIÓN DEL DOCUMENTO
\def\documenttitle {Estudio de caso: Poder, Estado y Derecho}
\def\documentsubtitle {Actividad 4 - Filosofía del Derecho}
\def\documentsubject {Semana 6 - Estudio de caso}

\def\documentauthor {Martín Jonathan de la Cruz}
\def\coursename {Filosofía del Derecho}
\def\coursecode {FDE30}

\def\universityname {Universidad Abierta y a Distancia de México}
\def\universityfaculty {Licenciatura en Derecho}
\def\universitydepartment {Filosofía del Derecho}
\def\universitydepartmentimage {departamentos/UnADM}
\def\universitydepartmentimagecfg {height=1.57cm}
\def\universitylocation {Roma Norte, Ciudad de México}

\def\authortable {
	\begin{tabular}{ll}
		\textbf{Alumno:}
		& \begin{tabular}[t]{l}
			Martín Jonathan de la Cruz \\
		\end{tabular} \\
		Matrícula: & ES2611202040 \\
		
		\textit{Profesora:}
		& \begin{tabular}[t]{l}
			Mtra. Reyna Patricia \\ González Rodríguez
		\end{tabular} \\
		& \\
		\multicolumn{2}{l}{Fecha de realización: 17 de mayo de 2026} \\
		\multicolumn{2}{l}{\universitylocation}
	\end{tabular}
}

% IMPORTACIÓN DEL TEMPLATE
\input{template}

% FORMATO SOLICITADO
\usepackage{helvet}
\renewcommand{\familydefault}{\sfdefault}
\usepackage{tikz}
\usepackage{longtable}
\usepackage{booktabs}
\usepackage{array}
\usepackage{ragged2e}

% CITAS AUTOR-AÑO
\setcitestyle{authoryear,open={(},close={)}}

% MARCA DE AGUA EN PORTADA
\def\coverwatermarkenabled {true}
\def\coverwatermarkimage {img/departamentos/UnADM.pdf}
\def\coverwatermarkopacity {0.16}
\def\coverwatermarkwidth {0.72\paperwidth}
\def\coverwatermarkxshift {0cm}
\def\coverwatermarkyshift {-0.35cm}

\newcommand{\insertcoverwatermark}{%
	\ifthenelse{\equal{\coverwatermarkenabled}{true}}{%
		\AddToShipoutPictureBG*{%
			\begin{tikzpicture}[remember picture,overlay]
				\node[
				opacity=\coverwatermarkopacity,
				inner sep=0pt,
				xshift=\coverwatermarkxshift,
				yshift=\coverwatermarkyshift
				] at (current page.center) {\includegraphics[width=\coverwatermarkwidth]{\coverwatermarkimage}};
			\end{tikzpicture}%
		}%
	}{}%
}

\newcolumntype{P}[1]{>{\RaggedRight\arraybackslash}p{#1}}

\begin{document}
	
	% PORTADA
	\insertcoverwatermark
	\templatePortrait
	
	% CONFIGURACIÓN DE PÁGINA Y ENCABEZADOS
	\templatePagecfg
	\onehalfspacing
	
	\begin{abstractd}
		\textbf{Objetivo:} Analizar la relación entre poder, Estado y derecho a partir del estudio de caso sobre el Estado mexicano en la Constitución Política de los Estados Unidos Mexicanos, identificando conceptos clave, artículos constitucionales y ejemplos periodísticos que muestran su operación en la práctica jurídica.
		
		\textbf{Postura de análisis:} El Estado de Derecho no se sostiene solo porque existan normas; se sostiene cuando el poder político, social y jurídico queda sometido a límites constitucionales, a división de funciones y a obligaciones reales frente a las personas.
	\end{abstractd}
	
	\templateIndex
	\templateFinalcfg
	
	\section{Introducción}
	
	El Estado mexicano no puede entenderse únicamente como gobierno o administración pública; funciona como una estructura de poder organizada jurídicamente. En esa estructura se cruzan territorio, población, soberanía, instituciones, normas, derechos humanos y mecanismos de control. Si uno de esos elementos falla, el derecho deja de operar como límite y empieza a funcionar solo como instrumento de mando.
	
	Desde la Filosofía del Derecho, el punto central no es repetir artículos constitucionales, sino identificar qué función cumplen dentro del sistema. El poder necesita organización para no convertirse en arbitrariedad; el Estado necesita derecho para justificar su actuación; y el derecho necesita instituciones capaces de aplicarlo sin romper la dignidad humana, la división de poderes ni la supremacía constitucional. En ese sentido, el estudio de caso permite observar cómo los artículos 1, 39, 40, 41, 42, 49 y 133 de la Constitución forman una arquitectura mínima del Estado mexicano \citep{cpeum2026, estudioCaso2026}.
	
	La posición que guía este trabajo es que el Estado de Derecho se juega en la práctica: en elecciones, fronteras, decisiones judiciales, reformas constitucionales, derechos de las víctimas y controles entre poderes. Por eso el análisis integra conceptos, artículos y notas periodísticas, para conectar la teoría con situaciones reales del sistema jurídico mexicano.
	
	\section{Definición de conceptos}
	
	\subsection{Concepto de Estado}
	
	El Estado es una organización política soberana que ejerce autoridad sobre una población asentada en un territorio determinado. No es solo el gobierno de turno; es el sistema institucional que produce, aplica y hace cumplir normas dentro de una comunidad. Paredes Lovón define el Estado a partir de una comunidad o población, un territorio, una organización política y una soberanía territorial \citep{paredes2020}. En términos prácticos, el Estado es el marco que convierte el poder disperso de la sociedad en poder institucional, y el derecho es el lenguaje normativo que vuelve exigibles sus decisiones \citep{garcia2002}.
	
	\subsection{Elementos del Estado}
	
	Los elementos básicos del Estado son población, territorio, gobierno y soberanía. La población representa el componente humano; el territorio fija el espacio donde se ejerce la autoridad; el gobierno organiza las instituciones que toman decisiones; y la soberanía expresa la capacidad de autodeterminación del Estado. Estos elementos no operan aislados: forman un sistema. Sin población no hay comunidad política; sin territorio no hay ámbito de validez; sin gobierno no hay dirección institucional; y sin soberanía no hay poder propio frente a otros poderes \citep{paredes2020}.
	
	\subsection{Formas de Estado}
	
	Las formas de Estado describen la manera en que se organiza el poder político. En una lectura básica, puede hablarse de monarquía y república; pero, en el análisis constitucional mexicano, lo relevante es la forma republicana, democrática, laica y federal. El artículo 40 constitucional establece que México se constituye como una República representativa, democrática, laica y federal \citep{cpeum2026}. Esto significa que el poder no pertenece a una persona ni a una dinastía, sino que se organiza mediante representación, elecciones, entidades federativas y límites jurídicos.
	
	\subsection{Sistema de gobierno}
	
	El sistema de gobierno indica cómo se ejerce el poder, cómo se eligen autoridades y cómo se distribuyen responsabilidades. Puede presentarse de forma democrática o autoritaria, según el origen y los límites del poder. En México, el sistema se sostiene formalmente en la soberanía popular, la representación política y la renovación periódica de autoridades. El artículo 41 constitucional conecta este punto con el ejercicio de la soberanía a través de los Poderes de la Unión y de los estados \citep{cpeum2026}. La consecuencia es clara: el gobierno no se legitima solo por mandar, sino por provenir de reglas constitucionales y procesos democráticos.
	
	\subsection{División de poderes}
	
	La división de poderes es un mecanismo de control institucional. Su función es evitar que una sola persona o corporación concentre el poder legislativo, ejecutivo y judicial. El artículo 49 constitucional divide el Supremo Poder de la Federación para su ejercicio en Legislativo, Ejecutivo y Judicial \citep{cpeum2026}. Esta división no elimina los conflictos entre poderes, pero crea un diseño para procesarlos jurídicamente. Sin división de poderes, el Estado puede conservar apariencia legal, pero pierde contrapesos reales.
	
	\subsection{Estado de Derecho}
	
	El Estado de Derecho significa que el poder público está sometido a normas, procedimientos, derechos y controles. No basta con que exista autoridad; esa autoridad debe actuar conforme a la Constitución, respetar derechos humanos y justificar sus decisiones. Kelsen advierte que el Estado puede entenderse como un orden jurídico, pero esa idea exige observar cómo se produce y se valida el derecho dentro de una estructura normativa \citep{kelsen1982}. Desde una postura práctica, Estado de Derecho implica legalidad, control del poder, seguridad jurídica, división institucional y protección efectiva de las personas; por eso la teoría de los derechos no puede separarse del diseño institucional que limita a la autoridad \citep{ansutegui2007}.
	
	\subsection{Poder político, social y jurídico}
	
	El poder político es la capacidad de dirigir decisiones públicas mediante instituciones de gobierno, elecciones, representación y administración. El poder social es la capacidad de influir desde fuera del aparato estatal: medios de comunicación, organizaciones civiles, grupos económicos, movimientos sociales o comunidades. El poder jurídico es la facultad reconocida por normas para crear, interpretar, aplicar o ejecutar derecho. La diferencia es importante: el poder político gobierna, el poder social presiona o influye, y el poder jurídico da forma obligatoria a decisiones mediante normas, sentencias, actos administrativos o procedimientos \citep{paredes2020,kelsen1982}.
	
	\section{Análisis del estudio de caso}
	
	El estudio de caso ubica al Estado mexicano dentro de la Constitución y permite observar cómo sus elementos no son conceptos sueltos, sino partes de un mismo sistema jurídico. La siguiente tabla relaciona artículos constitucionales con conceptos de poder, Estado y derecho, y agrega una nota periodística que permite ver el concepto en operación práctica.
	
	\begingroup
	\small
	\setlength{\tabcolsep}{4pt}
	\renewcommand{\arraystretch}{1.28}
	\begin{longtable}{@{}P{0.18\linewidth}P{0.14\linewidth}P{0.62\linewidth}@{}}
		\caption{Análisis constitucional del Estado mexicano y ejemplos periodísticos}\label{tab:analisis-caso}\\
		\toprule
		\textbf{Elemento} & \textbf{Artículo} & \textbf{Explicación y ejemplo} \\
		\midrule
		\endfirsthead
		\toprule
		\textbf{Elemento} & \textbf{Artículo} & \textbf{Explicación y ejemplo} \\
		\midrule
		\endhead
		\midrule
		\multicolumn{3}{r}{Continúa en la siguiente página} \\
		\endfoot
		\bottomrule
		\endlastfoot
		
		\textbf{Soberanía nacional} &
		\textbf{39 CPEUM} &
		La soberanía reside esencial y originariamente en el pueblo. Esto significa que el poder público no nace de la voluntad personal de una autoridad, sino de una base política superior: la comunidad nacional. La elección presidencial de 2024, reportada por Reuters al informar el triunfo electoral de Claudia Sheinbaum, muestra cómo la soberanía popular se expresa mediante el voto y legitima el acceso al poder político \citep{cpeum2026,reuters2024sheinbaum}. \\
		\midrule
		
		\textbf{Forma de Estado} &
		\textbf{40 CPEUM} &
		México se define como República representativa, democrática, laica y federal. Esto implica que el poder se organiza territorialmente entre Federación y entidades federativas. La discusión legislativa sobre la llamada supremacía constitucional, que requería la aprobación de congresos estatales, permite observar que el federalismo no es decorativo: los estados participan en decisiones constitucionales de alcance nacional \citep{cpeum2026,elfinanciero2024blindaje}. \\
		\midrule
		
		\textbf{Sistema de gobierno y ejercicio del poder} &
		\textbf{41 CPEUM} &
		El pueblo ejerce su soberanía por medio de los Poderes de la Unión y de los estados. En la práctica, esto exige autoridades electorales, reglas de competencia y procedimientos para renovar cargos públicos. La nota de \textit{El País} sobre la solicitud del INE para aplazar la elección judicial a 2028 refleja la dimensión operativa del sistema democrático: no basta convocar elecciones, también se requiere capacidad institucional para organizarlas de forma viable \citep{cpeum2026,elpais2026ine}. \\
		\midrule
		
		\textbf{Territorio} &
		\textbf{42 CPEUM} &
		El territorio nacional delimita el espacio donde el Estado ejerce soberanía. Esto incluye porciones continentales, islas, aguas, espacio aéreo y demás elementos reconocidos constitucionalmente. La nota de AS México sobre los ríos y cruces fronterizos entre México y Guatemala muestra cómo el territorio se traduce en fronteras, cruces, control estatal y relaciones internacionales concretas \citep{cpeum2026,trejo2024guatemala}. \\
		\midrule
		
		\textbf{División de poderes} &
		\textbf{49 CPEUM} &
		La división de poderes impide concentrar funciones legislativas, ejecutivas y judiciales en una sola autoridad. La reforma judicial mexicana ha generado debate precisamente porque toca la estructura del Poder Judicial y su relación con los otros poderes. La cobertura de Associated Press sobre el rechazo de la Suprema Corte a limitar la reforma judicial permite ver una tensión real: quién controla al poder y bajo qué reglas constitucionales \citep{cpeum2026,ap2024judicial}. \\
		\midrule
		
		\textbf{Derechos humanos} &
		\textbf{1 CPEUM} &
		El artículo 1 obliga a todas las autoridades a promover, respetar, proteger y garantizar derechos humanos. Esto convierte a la persona en límite del poder estatal. La nota de \textit{El País} sobre la negativa de amparo en el caso Sasha Sokol muestra cómo el acceso a la justicia, la reparación del daño y la protección de víctimas no son temas abstractos: son formas concretas de exigir que el Estado no abandone la dignidad humana \citep{cpeum2026,elpais2025sokol}. \\
		\midrule
		
		\textbf{Supremacía constitucional} &
		\textbf{133 CPEUM} &
		La Constitución es la Ley Suprema de la Unión, junto con las leyes federales y tratados celebrados conforme a ella. Este principio ordena la jerarquía normativa y obliga a jueces y autoridades a resolver conforme al bloque constitucional aplicable. La discusión periodística sobre la reforma de supremacía constitucional muestra que este principio no solo organiza normas: también define qué tan revisables son las reformas y qué margen conserva el Poder Judicial para proteger derechos \citep{cpeum2026,elfinanciero2024supremacia}. \\
		
	\end{longtable}
	\endgroup
	
	\section{Conclusiones}
	
	Hablar de Estado de Derecho exige conectar conceptos que a veces se estudian separados. Estado, gobierno, soberanía, división de poderes, derechos humanos y supremacía constitucional no son definiciones de memoria; son piezas de control. La población da base humana al Estado, el territorio delimita su campo de acción, la soberanía justifica el origen del poder, el gobierno lo organiza, la división de poderes lo limita y los derechos humanos le ponen una frontera ética y jurídica. Si el poder no encuentra esos límites, puede seguir produciendo actos legales en apariencia, pero se debilita la legitimidad del sistema.
	
	La conclusión de fondo es que el Estado de Derecho no depende únicamente de tener una Constitución escrita. Depende de que esa Constitución funcione como regla real frente al poder político, social y jurídico. Cuando una elección se organiza, una frontera se administra, una víctima busca reparación, una reforma constitucional se debate o un tribunal controla decisiones públicas, se está probando si el derecho limita al poder o si el poder está usando al derecho como simple herramienta. Mi postura es que la práctica profesional jurídica debe partir de esa vigilancia: no basta saber qué dice la norma, hay que identificar qué poder activa, qué límite impone y qué consecuencia produce para las personas.
	
	En el caso mexicano, los artículos analizados muestran un diseño constitucional robusto, pero también frágil si no se aplica con criterio. La soberanía popular puede justificar cambios profundos, pero no debe borrar controles; el federalismo distribuye poder, pero requiere coordinación; la supremacía constitucional ordena el sistema, pero no debe usarse para cerrar toda discusión sobre derechos; y la división de poderes solo sirve si cada poder conserva capacidad real de freno. Por eso, el Estado de Derecho es una estructura viva: se sostiene cada vez que una autoridad acepta que su poder no es absoluto y que la Constitución no es un adorno, sino el límite operativo del sistema.
	
	\bibliography{filosofia-del-derecho}
	
\end{document}
